% Standalone resolution manuscript generated from solution.md.
% Compile with XeLaTeX; author and review metadata are maintained in Markdown.
\documentclass[10pt,a4paper]{article}
\usepackage[margin=24mm,headheight=15pt]{geometry}
\usepackage{fontspec}
\setmainfont{DejaVuSerif.ttf}[BoldFont=DejaVuSerif-Bold.ttf,ItalicFont=DejaVuSerif-Italic.ttf,BoldItalicFont=DejaVuSerif-BoldItalic.ttf]
\setsansfont{DejaVuSans.ttf}[BoldFont=DejaVuSans-Bold.ttf,ItalicFont=DejaVuSans-Oblique.ttf]
\setmonofont{DejaVuSansMono.ttf}
\usepackage{amsmath,amssymb,unicode-math}
\setmathfont{latinmodern-math.otf}
\usepackage{xcolor,microtype,parskip,needspace,fancyhdr}
\usepackage{bookmark}
\usepackage{xurl}
\hypersetup{colorlinks=true,urlcolor=blue,linkcolor=blue,pdftitle={SP-12:
The chromatic nullity bound from Hall's theorem},pdfauthor={George
Stepaniants},pdflang={en-GB}}
\urlstyle{same}
\setlength{\emergencystretch}{3em}
\providecommand{\tightlist}{\setlength{\itemsep}{0pt}\setlength{\parskip}{0pt}}
\providecommand{\pandocbounded}[1]{#1}
\setcounter{secnumdepth}{0}
\allowdisplaybreaks[1]
\widowpenalty=10000
\clubpenalty=10000
\pagestyle{fancy}
\fancyhf{}
\fancyhead[L]{\small\sffamily APPLICATION NOTE / SP-12}
\fancyhead[R]{\small\sffamily 11 September 2026}
\fancyfoot[L]{\parbox[b]{0.9\textwidth}{\fontsize{8}{10}\selectfont Hall's
theorem; AI-assisted application note and independent automated-agent
review.}}
\fancyfoot[R]{\footnotesize\thepage}
\begin{document}
{\raggedright\Large\bfseries SP-12: The chromatic nullity bound from
Hall's theorem\par}
\vspace{0.4em}
{\normalsize\bfseries George Stepaniants\par}
{\small Department of Computing and Mathematical Sciences, California
Institute of Technology, Pasadena, California, USA.\par}
\vspace{0.6em}
\pagestyle{plain}

\textbf{Literature-dependent resolution.} The all-graph theorem used
here is due to H. Tracy Hall. This application note was prepared with
substantial ChatGPT/Codex assistance. A separate
\href{https://github.com/ajt60gaibb/OpenProblemsInNLA/blob/main/references/stepaniants-sp11-sp12-2026-09-11/verification/SP-11-SP-12-independent-review.md}{independent
Codex-agent review} checked this deduction and the essential proof of
Hall's theorem and returned \textbf{PASS}. This is automated-agent
mathematical review, not external human peer review, formal
verification, or a claim that the cited preprint has been refereed.

\subsection{Statement and external
input}\label{statement-and-external-input}

Let \(G\) be a finite simple undirected graph on \(n\ge1\) vertices. Let
\(\mathcal S(G)\) be the real symmetric \(n\times n\) matrices with
off-diagonal nonzero entries exactly at the edges of \(G\) and
unrestricted diagonal. A matrix \(A\in\mathcal S(G)\) has the strong
Arnold property (SAP) when the only real symmetric \(X\) satisfying

\[
AX=0,\qquad A\circ X=0,\qquad I_n\circ X=0
\]

is \(X=0\). Here \(\circ\) is entrywise product. Define

\[
\nu(G)=\max\{\operatorname{nullity}A:
A\in\mathcal S(G),\ A\succeq0,\ A\text{ has SAP}\}.
\]

The target is \(\nu(G)\ge\chi(G)-1\), where \(\chi(G)\) is the chromatic
number. The principal external input is Hall's \textbf{Corollary 3.22},
based on \textbf{Theorem 3.20}:

\[\nu(F)\ge\delta(F)\qquad\text{for every finite simple graph }F.\tag{H}\]

The source {[}H{]} is the preprint arXiv:2601.01211v1, submitted 3
January 2026. Its current record still listed only v1 when checked on 11
September 2026. This note supplies the remaining implication explicitly.
The independent review checked the essential proof of (H); it also
documents nonessential wording issues, including ``positive definite''
in the abstract where the definitions and construction use positive
semidefinite matrices.

\Needspace{7\baselineskip}

\subsection{Lemma: induced-subgraph
monotonicity}\label{lemma-induced-subgraph-monotonicity}

If \(H\) is an induced subgraph of \(G\), then \(\nu(H)\le\nu(G)\).

\subsubsection{Proof for adding one
vertex}\label{proof-for-adding-one-vertex}

It suffices to treat \(G\) obtained by adding one vertex \(v\) to \(H\).
Let \(H\) have \(h\) vertices. Choose a PSD matrix \(A\in\mathcal S(H)\)
with SAP and rank \(r\). We construct a PSD SAP matrix in
\(\mathcal S(G)\) with rank \(r+1\), preserving nullity.

Let \(\mathcal M_r\) be the smooth manifold of \(h\times h\) PSD
symmetric matrices of rank \(r\) near \(A\). Let \(P\) send a symmetric
matrix to its entries at unordered nonedges of \(H\). We first show that
the differential of \(P|_{\mathcal M_r}\) is surjective at \(A\).

In an orthonormal basis adapted to the range and kernel of \(A\),
tangent matrices have the form

\[
\begin{pmatrix}S&T\\T^T&0\end{pmatrix},
\]

where \(S\) is an arbitrary symmetric \(r\times r\) matrix and \(T\) is
arbitrary. Its orthogonal complement in the space of symmetric matrices
is exactly the symmetric matrices \(X\) with \(AX=0\). If the
differential of \(P\) were not onto, its adjoint would give a nonzero
such normal matrix supported only on off-diagonal nonedges of \(H\).
This contradicts SAP. The differential is therefore surjective.

Write \(z_i=1\) when vertex \(i\) is adjacent to \(v\) in \(G\) and
\(z_i=0\) otherwise, and set \(b=\epsilon z\). The submersion theorem
gives matrices \(D_\epsilon\in\mathcal M_r\), converging to \(A\) as
\(\epsilon\to0\), with

\[
(D_\epsilon)_{ij}=-b_i b_j
\]

at each nonedge \(\{i,j\}\) of \(H\). Define

\[
B_\epsilon=\begin{pmatrix}1&b^T\\b&D_\epsilon+bb^T\end{pmatrix}.
\]

Its Schur complement is \(D_\epsilon\), so \(B_\epsilon\) is PSD of rank
\(r+1\). Old nonedges cancel exactly; old edges remain nonzero for
sufficiently small \(\epsilon\) because they converge to those of \(A\).
For nonzero \(\epsilon\), new edges are exactly the nonzero coordinates
of \(b\). Thus \(B_\epsilon\in\mathcal S(G)\).

To check SAP, let \(Z_G\) be the fixed vector space of symmetric
matrices with zero diagonal and zero entries on every edge of \(G\). At

\[B_0=\operatorname{diag}(1,A),\]

the map \(X\mapsto B_0X\) on \(Z_G\) is injective. Its first row forces
the first row of \(X\) to vanish; symmetry then forces its first column
to vanish. The remaining block \(X_H\) satisfies \(AX_H=0\) and is
supported on off-diagonal nonedges of \(H\), so SAP for \(A\) forces
\(X_H=0\).

Injectivity of a linear map between finite-dimensional spaces is open in
its matrix entries: a nonzero maximal-column-rank minor remains nonzero
under sufficiently small perturbations. Hence \(X\mapsto B_\epsilon X\)
stays injective on \(Z_G\) for small \(\epsilon\). Since \(B_\epsilon\)
has exactly the pattern \(G\), this is its SAP.

Nullity has been preserved:

\[
\operatorname{nullity}B_\epsilon=(h+1)-(r+1)=h-r.
\]

Choose \(A\) with nullity \(\nu(H)\) and add the omitted vertices one at
a time. This proves the lemma. The rank-zero case is included:
\(\mathcal M_0=\{0\}\) is a zero-dimensional manifold, and surjectivity
is forced by SAP as above. No connectedness assumption is used.
\(\square\)

\Needspace{7\baselineskip}

\subsection{The chromatic bound: affirmative answer to
SP-12}\label{the-chromatic-bound-affirmative-answer-to-sp-12}

Set \(k=\chi(G)\). Choose an induced subgraph \(H\) minimal in vertex
count subject to \(\chi(H)=k\). For \(k\ge2\), deleting any vertex of
\(H\) leaves a graph colorable with at most \(k-1\) colors. If a vertex
had degree at most \(k-2\), such a coloring after its deletion would
leave an unused color among its neighbors and would extend to \(H\), a
contradiction. Therefore \(\delta(H)\ge k-1\).

Combine the lemma, Hall's theorem (H), and this degree inequality:

\[
\nu(G)\ge\nu(H)\ge\delta(H)\ge k-1=\chi(G)-1.
\]

For \(k=1\), the target \(\nu(G)\ge0\) is immediate. This proves the
exact all-graph assertion without assuming Hadwiger's conjecture.
\(\square\)

\subsection{Scope and attribution}\label{scope-and-attribution}

Hall's all-graph theorem is the essential external input. George
Stepaniants is the author of this explanatory application note, not the
discoverer of Hall's theorem. No novelty is asserted for
induced-subgraph monotonicity or the chromatic implication, which the
original problem source already identifies as following from Strong
Delta. The note does not provide a new independent proof of the central
existence theorem.

The
\href{https://github.com/ajt60gaibb/OpenProblemsInNLA/blob/main/references/stepaniants-sp11-sp12-2026-09-11/README.md}{submission
record} preserves the reconstructed source, independent review and dated
public-network audit. Unsupported earlier claims about unavailable
artifacts and graph-atlas computations remain withdrawn.

\subsection{References}\label{references}

{[}H{]} H. Tracy Hall, \emph{The Delta Theorem: A dimension bound for
faithful orthogonal graph representations}, arXiv:2601.01211v1, 3
January 2026, Theorem 3.20 and Corollary 3.22.
\href{https://arxiv.org/html/2601.01211v1}{Versioned primary text} ·
\href{https://arxiv.org/abs/2601.01211}{Current arXiv record}.

{[}R12{]} \emph{Open Problems in Numerical Linear Algebra},
\href{https://github.com/ajt60gaibb/OpenProblemsInNLA/blob/main/eigenvalues-and-inverse-problems/SP-12/README.md}{SP-12:
original canonical target and references}.
\end{document}
